% GENERATED FILE. Edit canonical lessons, then run npm run generate:books.
% canonical-source-hash: fnv1a64:313bafeb9c262e3d
% canonical-lessons: KA-C253-udaharane, KA-C253-ghatane, KA-C253-nenapu, KA-C253-jivana, KA-C253-niyama

\chapter{Example, Event, Memory, Life, Rule}
\label{ch:ka-example-event-memory-life-rule}
\hlchaptermodality{253}

\begin{chapteropening}
\textbf{By the end of this chapter:} \emph{I can say an example, an event, a memory, life and a rule.}

Complete the last lesson of chapter 253: I can say an example, an event, a memory, life and a rule.
\end{chapteropening}

\section[\texorpdfstring{\kn{ಉದಾಹರಣೆ} (udāharaṇe)}{ಉದಾಹರಣೆ (udāharaṇe)}]{\kn{ಉದಾಹರಣೆ} (udāharaṇe) — an example}

\label{lesson:KA-C253-udaharane}



\begin{quote}
Before the new one: say the Kannada for along with, then the Kannada for together.
\end{quote}

\subsection*{You'll want to know: \kn{ಉದಾಹರಣೆ}}
\textbf{\kn{ಉದಾಹರಣೆ}} — \emph{udāharaṇe} — \textquotedblleft{}an example\textquotedblright{}.

\begin{grammarlens}[title={What you've built}]
One more word for this chapter.
\end{grammarlens}

\subsection*{Your turn}
\begin{itemize}
\raggedright
  \item \emph{Say it:} \emph{udāharaṇe}
  \item \emph{Say it:} \emph{udāharaṇe}, once more
  \item \emph{Say it:} say \emph{udāharaṇe}
  \item \emph{Recall:} say the Kannada for along with, then the Kannada for together, then say \emph{udāharaṇe} again
\end{itemize}

\begin{tcolorbox}[breakable,colback=teal!4,colframe=teal!35!black,title={Before you move on}]
What does \kn{ಉದಾಹರಣೆ} mean? (\textquotedblleft{}An example\textquotedblright{}.) Say it once more.
\end{tcolorbox}
\section[\texorpdfstring{\kn{ಘಟನೆ} (ghaṭane)}{ಘಟನೆ (ghaṭane)}]{\kn{ಘಟನೆ} (ghaṭane) — an event, an incident}

\label{lesson:KA-C253-ghatane}



\begin{quote}
Before the new one: say the Kannada for together, then the Kannada for an example.
\end{quote}

\subsection*{You'll want to know: \kn{ಘಟನೆ}}
\textbf{\kn{ಘಟನೆ}} — \emph{ghaṭane} — \textquotedblleft{}an event, an incident\textquotedblright{}.

\begin{grammarlens}[title={What you've built}]
One more word for this chapter.
\end{grammarlens}

\subsection*{Your turn}
\begin{itemize}
\raggedright
  \item \emph{Say it:} \emph{ghaṭane}
  \item \emph{Say it:} \emph{ghaṭane}, once more
  \item \emph{Say it:} say \emph{ghaṭane}
  \item \emph{Recall:} say the Kannada for together, then the Kannada for an example, then say \emph{ghaṭane} again
\end{itemize}

\begin{tcolorbox}[breakable,colback=teal!4,colframe=teal!35!black,title={Before you move on}]
What does \kn{ಘಟನೆ} mean? (\textquotedblleft{}An event, an incident\textquotedblright{}.) Say it once more.
\end{tcolorbox}
\section[\texorpdfstring{\kn{ನೆನಪು} (nenapu)}{ನೆನಪು (nenapu)}]{\kn{ನೆನಪು} (nenapu) — a memory}

\label{lesson:KA-C253-nenapu}



\begin{quote}
Before the new one: say the Kannada for an example, then the Kannada for an event.
\end{quote}

\subsection*{You'll want to know: \kn{ನೆನಪು}}
\textbf{\kn{ನೆನಪು}} — \emph{nenapu} — \textquotedblleft{}a memory\textquotedblright{}.

\begin{grammarlens}[title={What you've built}]
One more word for this chapter.
\end{grammarlens}

\subsection*{Your turn}
\begin{itemize}
\raggedright
  \item \emph{Say it:} \emph{nenapu}
  \item \emph{Say it:} \emph{nenapu}, once more
  \item \emph{Say it:} say \emph{nenapu}
  \item \emph{Recall:} say the Kannada for an example, then the Kannada for an event, then say \emph{nenapu} again
\end{itemize}

\begin{tcolorbox}[breakable,colback=teal!4,colframe=teal!35!black,title={Before you move on}]
What does \kn{ನೆನಪು} mean? (\textquotedblleft{}A memory\textquotedblright{}.) Say it once more.
\end{tcolorbox}
\section[\texorpdfstring{\kn{ಜೀವನ} (jīvana)}{ಜೀವನ (jīvana)}]{\kn{ಜೀವನ} (jīvana) — life}

\label{lesson:KA-C253-jivana}



\begin{quote}
Before the new one: say the Kannada for an event, then the Kannada for a memory.
\end{quote}

\subsection*{You'll want to know: \kn{ಜೀವನ}}
\textbf{\kn{ಜೀವನ}} — \emph{jīvana} — \textquotedblleft{}life\textquotedblright{}.

\begin{grammarlens}[title={What you've built}]
One more word for this chapter.
\end{grammarlens}

\subsection*{Your turn}
\begin{itemize}
\raggedright
  \item \emph{Say it:} \emph{jīvana}
  \item \emph{Say it:} \emph{jīvana}, once more
  \item \emph{Say it:} say \emph{jīvana}
  \item \emph{Recall:} say the Kannada for an event, then the Kannada for a memory, then say \emph{jīvana} again
\end{itemize}

\begin{tcolorbox}[breakable,colback=teal!4,colframe=teal!35!black,title={Before you move on}]
What does \kn{ಜೀವನ} mean? (\textquotedblleft{}Life\textquotedblright{}.) Say it once more.
\end{tcolorbox}
\section[\texorpdfstring{\kn{ನಿಯಮ} (niyama)}{ನಿಯಮ (niyama)}]{\kn{ನಿಯಮ} (niyama) — a rule}

\label{lesson:KA-C253-niyama}



\begin{quote}
Before the new one: say the Kannada for a memory, then the Kannada for life.
\end{quote}

\subsection*{You'll want to know: \kn{ನಿಯಮ}}
\textbf{\kn{ನಿಯಮ}} — \emph{niyama} — \textquotedblleft{}a rule\textquotedblright{}.

\begin{grammarlens}[title={What you've built}]
One more word for this chapter.
\end{grammarlens}

\subsection*{Your turn}
\begin{itemize}
\raggedright
  \item \emph{Say it:} \emph{niyama}
  \item \emph{Say it:} \emph{niyama}, once more
  \item \emph{Say it:} say \emph{niyama}
  \item \emph{Recall:} say the Kannada for a memory, then the Kannada for life, then say \emph{niyama} again
\end{itemize}

\begin{tcolorbox}[breakable,colback=teal!4,colframe=teal!35!black,title={Before you move on}]
What does \kn{ನಿಯಮ} mean? (\textquotedblleft{}A rule\textquotedblright{}.) Say it once more.
\end{tcolorbox}
